\begin{titlepage}
\centering
{\Large Notas de una clase técnica con pantalla compartida\par}
\vspace{1.0cm}
{\huge\bfseries Modelos fundacionales para robots y VLA: hoja de ruta de los artículos clásicos}\par
\vspace{0.5cm}
{\Large Zhang Xiaojun Podcast conversa con Chen Jianyu (陈建宇)\par}
\vspace{0.35cm}
{\large Elaboradas a partir del video con pantalla compartida, la transcripción, la estructura de capítulos, la descripción del video y los fotogramas clave de la pantalla compartida\par}
\vspace{0.3cm}
{\large 2025-04-07\par}
\vspace{0.85cm}
\includegraphics[width=0.88\textwidth,height=0.42\textheight,keepaspectratio]{cover.jpg}\par
\vfill
\begin{tcolorbox}[width=0.92\textwidth, colback=black!2!white, colframe=black!55, sharp corners]
\textbf{Título del video}: Análisis, artículo por artículo, de los modelos fundacionales para robots y los artículos clásicos de VLA (con versión de pantalla compartida): «el ser humano es el VLA más inteligente»\par
\textbf{Canal del video}: Zhang Xiaojun Podcast\par
\textbf{Duración del video}: 02:29:41\par
\textbf{Enlace del video}: \href{\videourl}{\nolinkurl{\videourl}}\par
\textbf{Nota sobre el material}: YouTube no ofrece subtítulos descargables; el texto se elaboró a partir de una transcripción por bloques con faster-whisper large-v3 (CPU, int8), la descripción del video, la estructura de capítulos y los fotogramas clave de la pantalla compartida. Este video contiene mucho contenido en pantalla compartida, por lo que el texto usa fotogramas reales de esa pantalla como eje visual.\par
\textbf{Página del proyecto}: \href{\repourl}{\nolinkurl{\repourl}}\par
\end{tcolorbox}
\end{titlepage}

\begingroup
\small
\setlength{\parskip}{0pt}
\renewcommand{\baselinestretch}{0.92}\selectfont
\tableofcontents
\endgroup
\newpage

